\cotefile{78}
\scanpage{117}

on trouve

\[
\begin{aligned}
\Phi_{2p}(T)&=\Phi_p(-T)=T^{p-1}-T^{p-\unsure{2}}\ldots+1\\
\deleted{\unsure{\text{on a}}}\quad\Psi_{2p}(U)&=\Psi_p(-U)(-1)^{\frac{p-1}{2}}
\end{aligned}
\]

en car. $p$, les \deleted{solutions} \inserted{zéros de} \deleted{sont} $\Phi_{2p}(T)$ sont égales à $-1$, donc celles de $\Psi_{2p}$ sont égales à $-2$, qui est bien l'invariant du pol. régulier à $2p$ côtés. Donc en \uline{tte car.} $\ne 2$, la condition \unsure{nec.|n.} et \unsure{suffisante} pour qu'\deleted{\unsure{on ait}} l'inv. $\alpha$ définisse un pol. régulier à $2p$ côtés est\unsure{.|:} $\unsure{\Psi|\Phi}_{2p}(\alpha)=0$.

L'anneau modulaire est donc

\[
k'_{2p}=\mathbb Z[\tfrac12][U]/\Psi_{2p}(U).
\]

Reste enfin le cas $n=4$ ($=2.p$ avec $p=2$). Ici $\Phi_4(T)=T^2+1$, $\Psi_4(U)=\underset{\inserted{=F_4(U)}}{U}$, En car. 2, le\deleted{\unsure{s}} \inserted{seul} zéro de $\Psi_4$ est $0=-2$, donc qui définit bien le carré en car. 2. Donc ici

\[
{\unsure{k|K}}'_4=\mathbb Z[U]/U\simeq\mathbb Z
\]

i.e. il existe (à un isom unique près) une réalisation géom régulière de $\Pi_4$ et une seule sur un anneau $k$ donné, et celle-ci est \deleted{\unsure{eq}} fidèle quand l'anneau est un corps \ldots

% archive sheet number: 116
